% GENERATED FILE. Edit canonical lessons, then run npm run generate:books.
% canonical-source-hash: fnv1a64:fbfb068e3594aa17
% canonical-lessons: HI-C468-dhairya, HI-C468-adhir, HI-C468-jhakki, HI-C468-jhirakna, HI-C468-jhallana

\chapter{Patience and Temper}
\label{ch:hi-patience-and-temper}
\hlchaptermodality{468}

\begin{chapteropening}
\textbf{By the end of this chapter:} \emph{I can say whether someone keeps their patience.}

Complete the last lesson of chapter 468: I can say whether someone keeps their patience.
\end{chapteropening}

\section[\texorpdfstring{\hi{धैर्य} (dhairya)}{धैर्य (dhairya)}]{\hi{धैर्य} (dhairya) — forbearance; patience}

\label{lesson:HI-C468-dhairya}



\begin{quote}
Before the new one: say the Hindi for a betrayal of trust, then the Hindi for crooked.
\end{quote}

\subsection*{You'll want to know: \hi{धैर्य}}
\textbf{\hi{धैर्य}} — \emph{dhairya} — \textquotedblleft{}forbearance; patience\textquotedblright{}.

\begin{quote}
\hi{बच्चों} \hi{को} \hi{पढ़ाने} \hi{के} \hi{लिए} \hi{बहुत} \hi{धैर्य} \hi{चाहिए।}
\end{quote}

\begin{quote}
\emph{baccoṁ ko paṛhāne ke liye bahut dhairya cāhiye.}
\end{quote}

\begin{quote}
Teaching children takes a great deal of patience.
\end{quote}

Masculine.

\begin{cousinweb}
\textbf{\hi{धैर्य}} (\emph{dhairya}) is a tatsama, borrowed from Sanskrit \emph{dhairya}, \textquotedblleft{}firmness, constancy, calmness, patience\textquotedblright{} (Monier-Williams). It is the formal sister of \emph{dhīraj}.
\end{cousinweb}

\begin{grammarlens}[title={What you've built}]
One more word for describing what you lived and felt.
\end{grammarlens}

\subsection*{Your turn}
\begin{itemize}
\raggedright
  \item \emph{Say it:} \emph{dhairya}
  \item \emph{Say it:} \emph{dhairya}, once more
  \item \emph{Say it:} \emph{baccoṁ ko paṛhāne ke liye bahut dhairya cāhiye.} — Teaching children takes a great deal of patience.
  \item \emph{Recall:} say the Hindi for a betrayal of trust, then the Hindi for crooked, then say \emph{dhairya} again
\end{itemize}

\begin{tcolorbox}[breakable,colback=teal!4,colframe=teal!35!black,title={Before you move on}]
What does \hi{धैर्य} mean? (\textquotedblleft{}Forbearance; patience\textquotedblright{}.) Say it once more.
\end{tcolorbox}
\section[\texorpdfstring{\hi{अधीर} (adhīr)}{अधीर (adhīr)}]{\hi{अधीर} (adhīr) — impatient}

\label{lesson:HI-C468-adhir}



\begin{quote}
Before the new one: say the Hindi for crooked, then the Hindi for forbearance.
\end{quote}

\subsection*{You'll want to know: \hi{अधीर}}
\textbf{\hi{अधीर}} — \emph{adhīr} — \textquotedblleft{}impatient\textquotedblright{}.

\begin{quote}
\hi{ट्रेन} \hi{देर} \hi{से} \hi{आई}, \hi{इसलिए} \hi{सब} \hi{यात्री} \hi{अधीर} \hi{थे।}
\end{quote}

\begin{quote}
\emph{ṭren der se āī, isliye sab yātrī adhīr the.}
\end{quote}

\begin{quote}
The train came late, so all the passengers were impatient.
\end{quote}

\begin{cousinweb}
\textbf{\hi{अधीर}} (\emph{adhīr}) is a tatsama, borrowed from Sanskrit \emph{adhīra}, \textquotedblleft{}not steady, excitable\textquotedblright{}, built on \emph{a-}, \textquotedblleft{}not\textquotedblright{}, and \emph{dhīra}, \textquotedblleft{}steady\textquotedblright{} (Monier-Williams). Hindi uses it for someone who cannot wait.
\end{cousinweb}

\begin{grammarlens}[title={What you've built}]
One more word for describing what you lived and felt.
\end{grammarlens}

\subsection*{Your turn}
\begin{itemize}
\raggedright
  \item \emph{Say it:} \emph{adhīr}
  \item \emph{Say it:} \emph{adhīr}, once more
  \item \emph{Say it:} \emph{ṭren der se āī, isliye sab yātrī adhīr the.} — The train came late, so all the passengers were impatient.
  \item \emph{Recall:} say the Hindi for crooked, then the Hindi for forbearance, then say \emph{adhīr} again
\end{itemize}

\begin{tcolorbox}[breakable,colback=teal!4,colframe=teal!35!black,title={Before you move on}]
What does \hi{अधीर} mean? (\textquotedblleft{}Impatient\textquotedblright{}.) Say it once more.
\end{tcolorbox}
\section[\texorpdfstring{\hi{झक्की} (jhakkī)}{झक्की (jhakkī)}]{\hi{झक्की} (jhakkī) — cranky; set in odd ways}

\label{lesson:HI-C468-jhakki}



\begin{quote}
Before the new one: say the Hindi for forbearance, then the Hindi for impatient.
\end{quote}

\subsection*{You'll want to know: \hi{झक्की}}
\textbf{\hi{झक्की}} — \emph{jhakkī} — \textquotedblleft{}cranky; set in odd ways\textquotedblright{}.

\begin{quote}
\hi{मेरे} \hi{नाना} \hi{थोड़े} \hi{झक्की} \hi{हैं}, \hi{हर} \hi{काम} \hi{अपने} \hi{तरीक़े} \hi{से} \hi{करते} \hi{हैं।}
\end{quote}

\begin{quote}
\emph{mere nānā thoṛe jhakkī haiṁ, har kām apne tarīqe se karte haiṁ.}
\end{quote}

\begin{quote}
My grandfather is a bit cranky; he does everything his own way.
\end{quote}

\begin{cousinweb}
\textbf{\hi{झक्की}} (\emph{jhakkī}) is a tadbhava: CDIAL lists Hindi \emph{jhakkī}, \textquotedblleft{}stubborn\textquotedblright{}, under \emph{*jhakk}, a form it reconstructs (CDIAL 5319). It is the stubbornness of odd, fixed habits.
\end{cousinweb}

\begin{grammarlens}[title={What you've built}]
One more word for describing what you lived and felt.
\end{grammarlens}

\subsection*{Your turn}
\begin{itemize}
\raggedright
  \item \emph{Say it:} \emph{jhakkī}
  \item \emph{Say it:} \emph{jhakkī}, once more
  \item \emph{Say it:} \emph{mere nānā thoṛe jhakkī haiṁ, har kām apne tarīqe se karte haiṁ.} — My grandfather is a bit cranky; he does everything his own way.
  \item \emph{Recall:} say the Hindi for forbearance, then the Hindi for impatient, then say \emph{jhakkī} again
\end{itemize}

\begin{tcolorbox}[breakable,colback=teal!4,colframe=teal!35!black,title={Before you move on}]
What does \hi{झक्की} mean? (\textquotedblleft{}Cranky; set in odd ways\textquotedblright{}.) Say it once more.
\end{tcolorbox}
\section[\texorpdfstring{\hi{झिड़कना} (jhiṛaknā)}{झिड़कना (jhiṛaknā)}]{\hi{झिड़कना} (jhiṛaknā) — to snap at; to rebuke}

\label{lesson:HI-C468-jhirakna}



\begin{quote}
Before the new one: say the Hindi for impatient, then the Hindi for cranky.
\end{quote}

\subsection*{You'll want to know: \hi{झिड़कना}}
\textbf{\hi{झिड़कना}} — \emph{jhiṛaknā} — \textquotedblleft{}to snap at; to rebuke\textquotedblright{}.

\begin{quote}
\hi{देर} \hi{से} \hi{आने} \hi{पर} \hi{मालिक} \hi{ने} \hi{नौकर} \hi{को} \hi{झिड़का।}
\end{quote}

\begin{quote}
\emph{der se āne par mālik ne naukar ko jhiṛkā.}
\end{quote}

\begin{quote}
The master snapped at the servant for coming late.
\end{quote}

Past, with \hi{ने} (he or she …): \hi{उसने} \hi{नौकर} \hi{को} \hi{झिड़का}, \emph{usne naukar ko jhiṛkā}.

\begin{cousinweb}
\textbf{\hi{झिड़कना}} (\emph{jhiṛaknā}) is a tadbhava: CDIAL lists Hindi \emph{jhiṛaknā}, \textquotedblleft{}to shake, scold\textquotedblright{}, under \emph{*jhaṭ}, a form it reconstructs (CDIAL 5327). It is a short, sharp telling-off, sharper than \emph{ḍā̃ṭnā}.
\end{cousinweb}

\begin{grammarlens}[title={What you've built}]
One more word for describing what you lived and felt.
\end{grammarlens}

\subsection*{Your turn}
\begin{itemize}
\raggedright
  \item \emph{Say it:} \emph{jhiṛaknā}
  \item \emph{Say it:} \emph{jhiṛaknā}, once more
  \item \emph{Say it:} \emph{der se āne par mālik ne naukar ko jhiṛkā.} — The master snapped at the servant for coming late.
  \item \emph{Recall:} say the Hindi for impatient, then the Hindi for cranky, then say \emph{jhiṛaknā} again
\end{itemize}

\begin{tcolorbox}[breakable,colback=teal!4,colframe=teal!35!black,title={Before you move on}]
What does \hi{झिड़कना} mean? (\textquotedblleft{}To snap at; to rebuke\textquotedblright{}.) Say it once more.
\end{tcolorbox}
\section[\texorpdfstring{\hi{झल्लाना} (jhallānā)}{झल्लाना (jhallānā)}]{\hi{झल्लाना} (jhallānā) — to snap irritably}

\label{lesson:HI-C468-jhallana}



\begin{quote}
Before the new one: say the Hindi for cranky, then the Hindi for to snap at.
\end{quote}

\subsection*{You'll want to know: \hi{झल्लाना}}
\textbf{\hi{झल्लाना}} — \emph{jhallānā} — \textquotedblleft{}to snap irritably\textquotedblright{}.

\begin{quote}
\hi{बार} \hi{बार} \hi{वही} \hi{सवाल} \hi{सुनकर} \hi{पिता} \hi{झल्ला} \hi{गए।}
\end{quote}

\begin{quote}
\emph{bār bār vahī savāl sunkar pitā jhallā gae.}
\end{quote}

\begin{quote}
Hearing the same question again and again, my father snapped.
\end{quote}

Past (he …): \hi{वह} \hi{झल्ला} \hi{गया}, \emph{vah jhallā gayā}.

\begin{cousinweb}
\textbf{\hi{झल्लाना}} (\emph{jhallānā}) is inherited: CDIAL lists Hindi \emph{jhallānā}, \textquotedblleft{}to burn\textquotedblright{}, beside \emph{jhallā}, \textquotedblleft{}hot\textquotedblright{}, under \emph{*jhal}, \textquotedblleft{}flash\textquotedblright{}, a form it reconstructs (CDIAL 5352). The sense \textquotedblleft{}to flare up in irritation\textquotedblright{} is a turn CDIAL does not give, though the heat behind it is plain.
\end{cousinweb}

\begin{grammarlens}[title={What you've built}]
One more word for describing what you lived and felt.
\end{grammarlens}

\subsection*{Your turn}
\begin{itemize}
\raggedright
  \item \emph{Say it:} \emph{jhallānā}
  \item \emph{Say it:} \emph{jhallānā}, once more
  \item \emph{Say it:} \emph{bār bār vahī savāl sunkar pitā jhallā gae.} — Hearing the same question again and again, my father snapped.
  \item \emph{Recall:} say the Hindi for cranky, then the Hindi for to snap at, then say \emph{jhallānā} again
\end{itemize}

\begin{tcolorbox}[breakable,colback=teal!4,colframe=teal!35!black,title={Before you move on}]
What does \hi{झल्लाना} mean? (\textquotedblleft{}To snap irritably\textquotedblright{}.) Say it once more.
\end{tcolorbox}
